% ભાગ: first-order-logic
% પ્રકરણ: models-theories
% વિભાગ: expressing-relations

\documentclass[../../../include/open-logic-section]{subfiles}

\begin{document}

\olfileid{fol}{mat}{exr}

\olsection{\article{structure}\printtoken{S}{structure}માં સંબંધો વ્યક્ત કરવા}

\begin{explain}
!!{formula}sનો એક મુખ્ય ઉપયોગ !!a{structure}~$\Struct M$માંના
ગુણધર્મો અને સંબંધોને~$\Struct M$ની ભાષા~$\Lang L$ના પાયાના
સંકેતોના પદોમાં વ્યક્ત કરવાનો છે. તેનો અર્થ નીચે મુજબ છે:
$\Struct M$નું !!{domain} વસ્તુઓનો ગણ છે. !!{constant}s,
!!{function}s અને !!{predicate}sનું~$\Struct M$માં અનુક્રમે
$\Domain M$ની વસ્તુઓ, $\Domain M$ પરના વિધેયો અને $\Domain M$ પરના
સંબંધો વડે અર્થઘટન થાય છે. દાખલા તરીકે, જો~$\Obj A^2_0$,~$\Lang L$માં
હોય, તો~$\Struct M$ તેને સંબંધ~$R=\Assign{{\Obj A^2_0}}{M}$ નિયુક્ત
કરે છે. પછી સૂત્ર~$\Atom{\Obj A^2_0}{\Obj v_1,\Obj v_2}$ આ જ સંબંધને
નીચેના અર્થમાં \emph{વ્યક્ત કરે છે}: જો ચલ-નિયુક્તિ~$s$,~$\Obj v_1$ને
$a\in\Domain{M}$ અને~$\Obj v_2$ને $b\in\Domain M$ પર મોકલે, તો
\[
Rab \text{\quad ત્યારે અને માત્ર ત્યારે જ્યારે\quad}
\Sat{M}{\Atom{\Obj A^2_0}{\Obj v_1, \Obj v_2}}[s].
\]
નોંધો કે અહીં ચલ-નિયુક્તિઓ સામેલ કરવી પડે છે: આપણે માત્ર ``$Rab$
ત્યારે અને માત્ર ત્યારે જ્યારે
$\Sat{M}{\Atom{\Obj A^2_0}{a,b}}$'' કહી શકતા નથી, કારણ કે~$a$ અને~$b$
આપણી ભાષાના સંકેતો નથી; તેઓ~$\Domain{M}$ના !!{element}s છે.

આપણી પાસે માત્ર આણ્વિક !!{formula}s નથી; તાર્કિક સંયોજકો અને પરિમાણકો
વડે તેમને જોડી પણ શકીએ છીએ. તેથી વધુ જટિલ !!{formula}s એવા બીજા
સંબંધો વ્યાખ્યાયિત કરી શકે છે જે સીધા~$\Struct M$માં સમાવેલા નથી.
આવું કેવી રીતે કરવું, અને ખાસ કરીને !!a{structure}માં કયા સંબંધોને
વ્યાખ્યાયિત કરી શકીએ, તેમાં આપણને રસ છે.
\end{explain}

\begin{defn}
$!A(\Obj v_1,\dots,\Obj v_n)$ એ~$\Lang L$નું એવું !!a{formula} હોય
જેમાં માત્ર $\Obj v_1$,~\dots,~$\Obj v_n$ મુક્ત આવે, અને~$\Struct M$ એ
~$\Lang L$ માટેની !!a{structure} હોય. કોઈ પણ એવી ચલ-નિયુક્તિ~$s$ માટે
કે $s(\Obj v_i)=a_i$ ($i=1,\dots,n$),
\[
Ra_1\dots a_n \text{\quad ત્યારે અને માત્ર ત્યારે જ્યારે\quad}
\Sat{M}{\Atom{!A}{\Obj v_1,\dots,\Obj v_n}}[s]
\]
હોય ત્યારે અને માત્ર ત્યારે
$!A(\Obj v_1,\dots,\Obj v_n)$ \emph{સંબંધ}
$R\subseteq\Domain M^n$ને \emph{વ્યક્ત કરે છે}.
\end{defn}

\begin{ex}
અંકગણિતના પ્રમાણભૂત નિદર્શ~$\Struct N$માં !!{formula}
$\Obj v_1<\Obj v_2\lor\eq[\Obj v_1][\Obj v_2]$,~$\Nat$ પરના
$\le$ સંબંધને વ્યક્ત કરે છે. !!{formula}
$\eq[\Obj v_2][\Obj v_1']$ ઉત્તરગામી સંબંધ, એટલે કે એવો સંબંધ
$R\subseteq\Nat^2$ વ્યક્ત કરે છે જેમાં $Rnm$ ત્યારે હોય જ્યારે~$m$,
~$n$નો ઉત્તરગામી હોય. સૂત્ર~$\eq[\Obj v_1][\Obj v_2']$ પૂર્વગામી
સંબંધને વ્યક્ત કરે છે. !!{formula}s
$\lexists[\Obj v_3][(\eq/[\Obj v_3][\Obj 0] \land
  \eq[\Obj v_2][(\Obj v_1 + \Obj v_3)])]$ અને
$\lexists[\Obj v_3][\eq[(\Obj v_1 + {\Obj v_3}')][\Obj v_2]]$ બંને
$\Obj <$ સંબંધને વ્યક્ત કરે છે. તેનો અર્થ એ છે કે અંકગણિતની ભાષામાં
વિધેય-સંકેત~$<$ ખરેખર અનાવશ્યક છે; તેને વ્યાખ્યાયિત કરી શકાય છે.
\end{ex}
\sourcecorrection{OLMAT-002}{મૂળની
\texttt{expressing-relations.tex}, પંક્તિ~64માં બીજા અસ્તિત્વવાચક
સૂત્રની સમાનતાની જમણી બાજુ માત્ર~\texttt{v\_2} છે, જ્યારે આ જ દાખલામાં
દરેક વસ્તુ-ભાષાના ચલને~\texttt{\string\Obj} વડે લખાયું છે. અહીં
વાક્યરચનાત્મક રીતે સુસંગત~$\Obj v_2$ પુનઃસ્થાપિત કર્યું છે.}

\begin{explain}
આ વિચાર માત્ર કોઈ ચોક્કસ !!{structure}s માટે જ રસપ્રદ નથી; જ્યારે પણ
આપણે અભિપ્રેત નિદર્શ કે નિદર્શોને વર્ણવવા કોઈ ભાષા વાપરીએ છીએ, એટલે
કે સિદ્ધાંતો વિચારીએ છીએ, ત્યારે તે સામાન્ય રીતે રસપ્રદ છે. આવા
સિદ્ધાંતોમાં ઘણી વાર પાયાના સંકેતો તરીકે માત્ર થોડાં !!{predicate}s
હોય છે, પરંતુ તેઓ જે વ્યાપકક્ષેત્રનું વર્ણન કરે છે તેમાં બીજા ઘણા
સંબંધો અગત્યની ભૂમિકા ભજવે છે. જો આ બીજા સંબંધોને ભાષાના પાયાના
!!{predicate}sનું અર્થઘટન કરતા સંબંધો વડે વ્યવસ્થિત રીતે વ્યક્ત કરી
શકાય, તો આપણે કહીએ કે તેમને ભાષામાં \emph{વ્યાખ્યાયિત} કરી શકીએ છીએ.
\end{explain}

\begin{prob}
$\Lang L_A$માં નીચેના સંબંધોને વ્યાખ્યાયિત કરતાં !!{formula}s શોધો:
\begin{enumerate}
\item $n$,~$i$ અને~$j$ની વચ્ચે છે;
\item $n$,~$m$ને નિઃશેષ ભાગે છે (એટલે કે~$m$,~$n$નો ગુણિત છે);
\item $n$ અવિભાજ્ય સંખ્યા છે (એટલે કે $1$ અને~$n$ સિવાય કોઈ સંખ્યા
  ~ $n$ને નિઃશેષ ભાગતી નથી).
\end{enumerate}
\end{prob}

\begin{prob}
ધારો કે સૂત્ર~$!A(\Obj v_1,\Obj v_2)$, !!a{structure}~$\Struct M$માં
સંબંધ~$R\subseteq\Domain M^2$ને વ્યક્ત કરે છે. નીચેના સંબંધો વ્યક્ત
કરતાં સૂત્રો શોધો:
\begin{enumerate}
\item $R$નો વ્યસ્ત~$R^{-1}$;
\item સાપેક્ષ ગુણાકાર~$R\mid R$.
\end{enumerate}
$R$નું પરંપરિત સંવરણ~$R^+$ વ્યક્ત કરવાની કોઈ રીત શોધી શકો છો?
\end{prob}

\begin{prob}
$\Lang{L}$ માત્ર દ્વિસ્થાની વિધેય-સંકેત~$<$ ધરાવતી ભાષા હોય (બીજાં
કોઈ !!{constant}s, !!{function}s કે !!{predicate}s નહિ---અલબત્ત,
$\eq$ સિવાય). $\Struct{N}$ એવી સંરચના હોય કે $\Domain{N}=\Nat$ અને
$\Assign{<}{N}=\Setabs{\tuple{n,m}}{n<m}$. નીચેનું સાબિત કરો:
\begin{enumerate}
\item $\{0\}$,~$\Struct{N}$માં વ્યાખ્યેય છે;
\item $\{1\}$,~$\Struct{N}$માં વ્યાખ્યેય છે;
\item $\{2\}$,~$\Struct{N}$માં વ્યાખ્યેય છે;
\item દરેક $n\in\Nat$ માટે ગણ~$\{n\}$,~$\Struct{N}$માં વ્યાખ્યેય છે;
\item $\Domain{N}$નો દરેક સાન્ત ઉપગણ~$\Struct{N}$માં વ્યાખ્યેય છે;
\item $\Domain{N}$નો દરેક સહસાન્ત ઉપગણ~$\Struct{N}$માં વ્યાખ્યેય છે
  (જ્યાં $X\subseteq\Nat$ ત્યારે અને માત્ર ત્યારે સહસાન્ત છે જ્યારે
  $\Nat\setminus X$ સાન્ત છે).
\end{enumerate}
\end{prob}


\end{document}
